\section{Discussion and Limitations}
\label{sec:discussion}
\paragraph{What the results show.} Cheap filters give only partial protection, how much depends on the attack, and adaptive attacks remove almost all of it. Checking how the pair's score and embedding change across several filters worked much better. It separated nearly all the gradient-based attacks we tried from clean images and harmless noise, including DeepFool and CW (never seen in training), smaller budgets and the adaptive attacks, at the cost of a modest number of false alarms. FGSM at $\epsilon{=}8$ is the clear exception. The face model itself is never retrained; the only parts we learn are small classifiers on an 11-number feature vector.

\paragraph{Why we do not claim the system is robust.} (i) \emph{Our adaptive attacks are not the worst case.} The detector-aware attack is one particular PGD variant (raise the similarity across filters and keep the embeddings steady). We never attacked the detector's decision directly, for example with a differentiable copy of the detector, or with more steps and restarts. A stronger attacker might get through. (ii) \emph{Flagging has a cost.} A flagged pair is either refused or sent for manual review. That happens to 5.3\% of clean pairs at our setting, and harmless noise is flagged a bit above the budget (8.3\%). (iii) \emph{Square.} The detector did not flag Square changes, and our Square and DeepFool attacks rarely succeed, so we cannot say how the system would do against stronger black-box or minimal-change attacks. (iv) \emph{Different budget.} At $\epsilon{=}8$ the detector misses 29\% of FGSM pairs, because it was trained at $\epsilon{=}4$.

\paragraph{Statistical caveats.} We used one seed, 300 attacked pairs (150 for adaptive), one dataset and one model. For a rate near 0.5, the 95\% interval is about $\pm$0.06 with 300 pairs and $\pm$0.08 with 150, so small gaps between filters should not be taken too seriously. Randomness comes from the PGD random start and from the classifiers. We used a center crop instead of face detection and alignment. We also did not look at differences between demographic groups (LFW is not balanced, and the detector might raise false alarms more often for some groups), and we did not check whether people in LFW were also in the model's training data.

\paragraph{Cost.} The detector needs six embeddings per image instead of one, so on our CPU it takes about \mscratio$\times$ the time of plain verification. That is cheap compared to adversarial training, but not ``free''.

\paragraph{Future work.} Attack the detector directly (surrogate gradients, more steps, restarts, EOT); try stronger black-box and minimal-change attacks; use more seeds, datasets and models with proper face alignment; tune the flag threshold for each real deployment; and check whether the flag rate is fair across groups.

\paragraph{Conclusion.} Simple image filters on their own are not a reliable defense for face verification once the attacker adapts to them. Checking whether a pair stays consistent across a few cheap filters, and combining that with the verification score through a Bayes rule, caught nearly all the gradient-based attacks in our tests (except FGSM at the larger budget). What is still missing is evidence that it holds up against a fully adaptive attacker.
